\begin{frame}{트리 개수를 늘리면: 정확도는 오르고 흔들림은 줄어든다}
  \begin{columns}[T]
    \begin{column}{0.52\textwidth}
      \centering
      \includegraphics[width=\linewidth]{ch07_2_n_estimators_curve.png}
      \vskip0.3em
      {\footnotesize test 정확도(선)와 표준편차(오차막대)를
      $n_{\text{estimators}}$(로그 스케일)의 함수로.}
    \end{column}
    \begin{column}{0.45\textwidth}
      $n_{\text{estimators}}$ 1$\to$200까지, 각 시드 10개 평균:
      \begin{enumerate}
        \item 정확도는 1$\to$10개 정도로 \textbf{빠르게 올랐다가} 거의
          \textbf{수평}(수확체감)
        \item 표준편차(오차막대)는 트리가 늘수록 \textbf{줄어든다} —
          ``전문가''가 많을수록 소수 의견에 흔들리지 않음
        \item 정확도가 \textbf{떨어지는 지점이 없다} — 트리를 더
          늘려도 \textbf{과적합이 안 온다}(7.3 GBDT와 대조) $\to$ 실전은
          정체 지점보다 여유를 두고 \textbf{100$\sim$300}이 안전
      \end{enumerate}
    \end{column}
  \end{columns}
\end{frame}
